\chapter{Discussion and Conclusions}
\label{sec:discussion}

\section{Discussion of the Results}

The chapter is organized around five findings, in the order in which the campaign produced them, with two engineering lessons attached to the ones they belong to.

The first finding is the one the thesis set out to test: an important part of the cost of declarative dynamic heuristics can be moved from grounding to search. On BSP, native \texttt{\#heuristic} directives grow as $\Theta(n^3)$ and exceed one million ground objects at $n=100$; the lazy variants keep two facts, and the difference is visible in the propositional instance size, in the peak memory, and in the timeout profile. The mechanism generalizes, and the full campaign measures the generalization. On PUP, whose heuristics multiply four unrolled value domains per ground pair, the ground Alpha simulation and the lazy variant both clear the entire published range under the 600-second/30~GB budget. The ground simulation, however, costs about twice the time and twice the memory of the lazy variant at every one of the ten common sizes, because the unrolled directives keep growing with the domains being unrolled. Both dominate the unguided baseline, which stalls by timeout at less than half their reach.

The two families therefore show the same mechanism from two sides. Where the heuristic grounding is the binding constraint, the lazy representation converts directly into a large cost advantage over the best ground heuristic, and into solved instances relative to the unguided baseline. Where the base program dominates, as at the top of the BSP range, it converts instead into an unchanged frontier reached with the search component reduced by three orders of magnitude; the residual wall there belongs to the domain encoding, not to the heuristic. The BSP frontier is unchanged rather than improved: the lazy variant, the clingo-like one and the heuristic-free baseline all stop within one size of each other, and the runs that decide the exact size finish within a few percent of the time limit. On that family the claim is about cost rather than reach, and the campaign's design --- structural counters alongside wall-clock times --- is what keeps the two apart.

PUP is the family where the claim is about reach as well, but the shape of that claim needs stating carefully, because it is bounded by the benchmark rather than by the method. Against the unguided and the clingo-semantics baselines the separation is a frontier: they stall at $N \le 100$ while the two Alpha-semantics variants solve every published instance. Between those two variants, \texttt{ga} and \texttt{la}, there is no separating instance, for the simple reason that \texttt{double-200} is the largest one the reference material contains and both solve it. What separates them there is cost at equal reach, and it is not small: $438.6$\,s and $24.4$\,GB against $224.0$\,s and $12.7$\,GB, on a $30$\,GB budget that \texttt{ga} is within twenty percent of exhausting. Extrapolating the two curves puts \texttt{ga}'s ceiling around $N \approx 215$ and \texttt{la}'s around $N \approx 280$ (Section~\ref{sec:pup}), so the measured factor of two in memory should be worth roughly a third in instance size; but that is an extrapolation, it is reported as one, and turning it into a measurement needs instances that do not exist yet.

The external comparison of Section~\ref{sec:alpha} puts a number on both halves of that sentence, and it is the most informative single experiment of the campaign because it is the only one whose yardstick is not clingo. Against Alpha in its dynamic-aggregate variant --- the system whose semantics this thesis reproduces, running the authors' own encodings --- the lazy propagator delivers the guidance and not the grounding. On BSP it reproduces the reference trajectory exactly, $n$ decisions and no backtracking, and then loses two orders of magnitude in wall time to a system that never builds the $52.8$ million rules of the base program. On PUP, where the base program is proportionally less dominant, it lands within a small factor of the reference in time and at a peak memory the two systems cannot be told apart on, while the ground simulation of the very same heuristic is nearly an order of magnitude slower than Alpha and uses almost twice its memory; Section~\ref{sec:alpha} gives the figures. Read as a whole, the ordering is \texttt{ga} $<$ \texttt{la} $<$ Alpha: the lazy representation recovers most of the distance, and grounding accounts for the remainder. That is the claim the thesis makes, now measured against an external reference rather than asserted.

The same experiment also protects the semantic claim from an obvious objection, namely that the Alpha reading only looks good because it is being compared with clingo baselines. Alpha \emph{without} its declarative heuristic stops at $n=40$ on BSP, eleven sizes before ground clingo without a heuristic, and solves no PUP instance at all. Lazy grounding and dynamic heuristics are therefore separable contributions and neither is sufficient on its own: the four combinations of (eager, lazy) grounding and (absent, dynamic) heuristic land at $n=150$, $n=160$, $n=40$ and $n=200$ respectively. The effect this thesis isolates inside clingo is the same effect that makes the reference system usable.

This result should not be read as a replacement for general lazy grounding. clingo still grounds the base program; the propagator acts only on the heuristic component, and Section~\ref{sec:alpha} measures exactly what that leaves on the table. Nor should it be read as a claim that laziness is free: the runtime work is real, it is paid at every decision, and the per-decision metrics exist precisely to keep that cost on the record. The honest formulation is a space--time trade-off with a crossover: below it, ground-and-solve remains competitive, because the explosion is manageable and the lazy machinery (in the Prolog case, an entire in-process logic-programming engine) carries a fixed overhead; beyond it, the constant-size heuristic component is the only representation that survives.

The second finding is semantic, and it is a finding about \emph{meaning}, not performance. The same declarative heuristic produces opposite outcomes depending on how partial information is read. Under the Alpha-like reading, the BSP heuristic is a perfect greedy strategy ($n$ choices, zero conflicts) and the HRP policy reproduces its intended level structure decision for decision. Under the clingo-like reading, the very same weights and bodies degenerate in two distinct ways, which Section~\ref{sec:clingo-semantics} documents together. On BSP and PUP the guards never open at decision time: the propagator reports zero applicable candidates per call, and the search is bit-identical to the baseline --- a correspondence that survives the change of domain encoding, since \texttt{lc\_co} likewise matches \texttt{gc\_noheur\_co} decision for decision. Call this inertness. On HRP the surviving low-level modifications steer the solver into a region almost eight hundred times larger at $N=14$ and five orders of magnitude larger at $N=16$, with both backends independently reaching the same order of magnitude of choices even though their exact counts differ. Call this misguidance.

Neither reading is \enquote{wrong}: they answer different questions, and the four-variant matrix exists to keep those questions separate. What the campaign shows is that the choice between them is not a tuning parameter but a precondition, because these encodings were written under one semantics and the other one cannot read their guards. A comparison that silently mixed the axes, for example by measuring \texttt{la} against \texttt{gc} and attributing the difference to laziness, would be methodologically invalid; the thesis takes some care never to do so. The clearest practical corollary is the one about the ground side: an inert heuristic is not a neutral heuristic if it has been materialized. On BSP the never-firing directives of \texttt{gc} cost two sizes of frontier against the baseline that carries no heuristic at all, $n=130$ against $n=150$, and on PUP they make \texttt{gc} fail before the baseline does. The lazy representation is therefore worth having even when the guidance is useless, which is not the argument the thesis set out to make and is arguably the more robust one.

The campaign adds a corollary about the ground simulation of the Alpha reading. \texttt{ga} reproduces the lazy guidance exactly while its weights remain representable, and stops working when they do not: on BSP that happens at $n=51$, and from $n=70$ the variant no longer finishes at all, making it the earliest-failing variant of the main matrix on BSP, nine sizes before the baseline. What stops it is not the price of the unrolling, which Section~\ref{sec:ga-weak} measures and finds negligible against the base program; it is the field the reconstructed aggregate value has to fit into. Inside a ground-and-solve pipeline the lazy propagator is therefore not merely the cheaper carrier of the Alpha semantics; beyond a certain instance size it is the only one that can carry it at all, because it ranks its targets internally and never has to encode the aggregate as a weight.

The \texttt{\_co} family sharpens that statement by supplying the case where the ceiling does no harm. \texttt{ga\_co} runs into it on exactly the same weights and still clears every benchmarked size, because what the saturation destroys is the ordering among candidates, and an encoding whose propagation already determines the partition never consults that ordering. The ceiling is thus a limit on the ranking rather than on the encoding, and it becomes visible in proportion to how much the run depends on the heuristic --- which is the same condition under which the heuristic is worth having in the first place.

The third finding generalizes that observation into a statement about the interface between the heuristic language and the solver that consumes it. clingo represents the weight of a \texttt{\#heuristic} modification in a bounded field and saturates anything larger, so a weight is not an arbitrary integer but a resource with a ceiling. Any encoding that packs information into the weight, whether an Alpha level through the flattening rule of Section~\ref{sec:flattening} or the value of an aggregate through the unrolling of \texttt{ga}, is spending that resource, and the two uses compete. Nothing warns the modeller when the ceiling is reached: the program still grounds, the run still terminates, the answer sets are still correct, and only the search degrades. The lazy propagator sidesteps the ceiling by construction, since it ranks targets internally and hands the solver a decision rather than a weight, and that is a design advantage worth stating separately from the grounding one.

The fourth finding concerns the translation layer between the two heuristic languages. Alpha's weight--priority pairs are global levels; clingo's priorities arbitrate on a single atom. Every encoding that expresses an ordering between different atom families must therefore flatten its levels into weights wherever clingo's reading applies: in the ground variants, because clasp is the consumer, and in the lazy clingo variants, because the propagator deliberately reproduces clingo's behaviour there. This is easy to state and easy to get wrong: an early revision of the lazy-clingo PUP and HRP encodings carried Alpha-style priorities that the clingo semantics does not rank globally, producing an unintended decision order that inverted the designed one (sensors before zones; cabinet-filling before reuse), and diverging from the Prolog twin encodings of the same variants. The final audit caught and corrected the discrepancy by applying the same flattening used in the ground encodings. The episode is instructive: the interaction between priority semantics and evaluation backend is exactly the kind of silent, non-crashing behavioural divergence that a benchmark pipeline cannot detect from exit codes, and it motivates the equivalence and wiring guards that the suite now includes.

The lesson was then turned into a structural decision rather than a matter of vigilance. As long as the priority field is used to carry levels in \emph{some} encodings, its meaning depends on the pair (semantics, backend): a global level under one reading, a purely local arbiter under another, and a global level again in the Prolog backend, whose selection rule sorts by (priority, weight) regardless of the semantics selector. Three readings of one syntactic field is a standing invitation to the very bug just described. Since flattening is order-preserving, the suite removes the ambiguity by construction: the levels always live in the weight, in every variant and in both backends. The native heuristic language went one step further and dropped the priority construct altogether, so that a directive written for it cannot express a level except through its weight, and the parser rejects the field rather than silently reinterpreting it. Two properties follow. First, the encodings of a pair differ by exactly one token, so the axis under study is the only thing that varies. Second, the change is behaviour-preserving with respect to the earlier revision: on every benchmark the flattened weights reproduce the previous total order exactly, because on each instance the level multiplier strictly dominates the intra-level weights, so the experimental results reported here are unaffected.

That decision has a cost, and Section~\ref{sec:dsl-limits} states it: a fragment with no priority field cannot express a genuinely two-dimensional rank whenever the intra-level weight range is not known before grounding. The two backends consequently do not cover the same language, the native one being a proper subset of the Prolog one, and the encodings of this suite were chosen to sit in the intersection so that the same experiment could run on both. That is a property of the benchmark selection and not of the mechanism, and it is the main thing a reader should not generalize from these results.

A related engineering lesson comes from the silent-fallback failure mode of the Prolog backend. A single missing environment variable degrades every lazy Prolog run to the no-heuristic baseline without any error, since the fallback evaluator simply fails on the Prolog-style rule bodies. The pipeline defends itself with positive evidence of activity (\mintinline{text}|decide_calls > 0|), cross-backend agreement checks, and a wrapper that forces the variable. Three of the episodes recorded in this chapter --- the saturating weight field, the reinterpreted priority annotation, and this one --- belong to the same class: the program grounds, the run terminates, the answer sets are correct, and only the guidance is silently gone. They argue for the same remedy, that a heuristic mechanism should check its own preconditions and prove it ran rather than merely not fail.

The fifth finding qualifies the first one, and it is the least comfortable of the set. On BSP the entire search advantage of the lazy variants rests on a domain encoding whose difficulty is self-inflicted. Restating the balance constraint compactly costs one line, and the resulting problem is solved by the unguided baseline in four milliseconds and sixty decisions at $n=200$ --- fewer decisions than any variant whose heuristic actually fires, and exactly as many as \texttt{lc\_co}, whose heuristic never does. The heuristic, in other words, is not merely unnecessary on the compact encoding; it is very slightly counterproductive, because it insists on deciding elements that propagation would have fixed. A reader is entitled to conclude from this that the BSP heuristic was never the interesting thing about BSP.

What survives the observation, and is strengthened by it, is the representation claim. Table~\ref{tab:co-family}, which reports that four-millisecond baseline, also shows the three ground variants spending between $26$ and $44$ seconds and up to $4.4$\,GB to express a heuristic that changes nothing, against twenty milliseconds and two facts for the lazy ones. With the base program reduced to six hundred rules there is no other explanation available for that spread: it is the heuristic component, measured on its own for the first time in the campaign. The practical reading is that the cost of \emph{carrying} a declarative heuristic and the value of \emph{having} one are independent quantities, and that the ground-and-solve pipeline couples them in the worst way, charging most for the heuristics it can least afford to unroll. That is an argument for the lazy representation which does not depend on the heuristic being good, and it is the one this thesis would keep if it had to keep only one.

Finally, the auxiliary-form experiment delimits the approach from within. \texttt{la\_aux} shows that laziness only pays if the \emph{entire} dynamic part of the heuristic, aggregate included, stays out of the grounder: materializing the heuristic body into an auxiliary relation reintroduces the very product the mechanism was built to avoid. This has a direct bearing on the expressiveness question above, because precompiling into auxiliary predicates is exactly the workaround by which the native fragment could accept a body it cannot express directly. The measurement says that this workaround is not a workaround: what the language cannot express, the encoding cannot cheaply supply on its behalf.

\section{Limitations}

The current implementation has explicit technical limitations. The linguistic ones are not restated here: Section~\ref{sec:dsl-limits} gives the native fragment in full along its five axes --- target, body, aggregate, weight and annotation --- and the discussion above has already drawn the consequence, that the two backends do not carry the same language. Two limitations of the implementation itself are not linguistic and belong in this list instead: the propagator is registered sequentially and has no synchronization layer for parallel solving, and the tie-break between targets of equal rank is deterministic but does not reproduce clingo's internal scores.

On the evaluation side, the campaign covers all three families on both backends with the corrected encodings, and both backends are phase-timed, but its scope needs stating precisely. Every cell of the design --- one variant, on one instance, on one backend --- is a single run, so the campaign measures no within-cell variance. Time-based claims are therefore reported as single measurements rather than averages. The separations this thesis rests on --- the factor of two between \texttt{ga} and \texttt{la} on PUP, the two orders of magnitude of \texttt{lc} on HRP, the three orders of magnitude on the BSP solving component --- are wide enough that run-to-run noise cannot account for them. The exception is the exact position of the BSP frontier, where the deciding runs finish within $6\%$ of the time limit; repeated runs per cell would settle those two sizes, and remain the most valuable extension of the protocol.

The benchmark sets bound the conclusions in three further ways. The HRP instances remain easy for clasp at every benchmarked size, so HRP supports conclusions about guidance quality and per-decision cost, not about frontiers. On BSP the bound is sharper and is discussed above: the difficulty that the heuristic overcomes is a property of the quadratic balance constraint, and the \texttt{\_co} family shows it dissolving under a one-line reformulation. Every BSP search result in this thesis should therefore be read as a measurement of guidance on a given encoding, not as evidence that the problem needs guidance. On PUP the claims rest on one instance per size in a single family, \texttt{double}, which ends at \texttt{double-200}: the non-monotone behaviour observed for \texttt{gc} and \texttt{lc} is a reminder that neighbouring sizes differ in hardness, per-size variance is not measured, and the comparison between \texttt{ga} and \texttt{la} is a comparison of cost on shared instances rather than of reach, because no available instance separates them. A further caveat concerns the Prolog backend specifically: its synchronization layer was rebuilt partway through the work, and the campaign reported here measures the current design throughout. The comparison against the previous one, quoted in Section~\ref{sec:per-decision-cost}, is therefore a comparison across two campaigns rather than within one, and while the effect is far too large to be noise --- a factor of $27$ on synchronization --- the per-family totals it is embedded in also carry the ordinary run-to-run drift discussed above.

The two backends also do not reproduce identical search trajectories on PUP and HRP, where the heuristic ranking admits ties that the C++ and Prolog tie-breaks resolve differently. The divergence reaches fourteen percent on \texttt{la} and is larger still on the unstable \texttt{lc} searches, which is why the cross-backend comparisons in those families are read qualitatively rather than as controlled measurements of the engine alone. The memory metric is end-to-end process RSS, which is the honest but blunt instrument discussed in Chapters~\ref{sec:implementation} and~\ref{sec:results}: it includes the Prolog engine for the lazy Prolog variants and cannot separate grounding memory from search memory.

The external comparison with Alpha inherits three limitations of its own. It carries different weight on different families: on BSP and PUP it measures the mechanism this thesis is about, whereas on HRP, whose reference heuristic contains no dynamic aggregates, it measures the cost of lazy grounding on a program that does not need it, and the divergence observed there between Alpha's two guided configurations is reported without being diagnosed. The two solvers' internal counters are not commensurable, so the comparison rests on the two \texttt{runlim} measures alone. And Alpha's peak RSS includes the JVM heap reserved through \texttt{-Xmx} rather than the heap in use, which makes its memory figures an upper bound, and the memory comparison on PUP a tie read conservatively rather than a measured equality.

Finally, the auxiliary clingo evaluator of the Prolog build re-grounds a program per decision and is a functional reference, not a performance-comparable backend, and the whitelist of static predicates together with the inference of the program constant $n$ in the Prolog backend are benchmark-specific conveniences that a production version should replace with a declarative interface.

\section{Implications}

The main implication is that declarative heuristics can be treated as an intermediate layer between modelling and solving. One does not have to choose between fully grounded \texttt{\#heuristic} directives and procedural heuristics hard-coded in the solver: a compact declarative surface, evaluated at runtime with access to the partial assignment, is a workable third point in the design space, and it can host more than one evaluation semantics behind one syntax. The two backends demonstrate two implementation strategies for that layer, a compiled, incremental one and an embedded general-purpose engine, with the expected trade-off between per-decision efficiency and expressive freedom of the rule bodies.

From an engineering perspective, the work also shows that clingo's propagator and heuristic interfaces are flexible enough to prototype non-trivial guidance mechanisms without touching the solver core: watched literals, incremental aggregate states, per-target modifier resolution, and a global decision ranking are all expressible against the public API of \texttt{clingo::Heuristic}.

\section{Future Work}
\label{sec:future-work}

\paragraph{Widening the native fragment.}
The single most valuable extension is the one Section~\ref{sec:dsl-limits} identifies: aggregates and body literals that can be joined to each other rather than only to the target tuple, together with comparison guards between bound values, integer division and modulo, and a principled treatment of non-numeric symbols through a stable internal mapping. The concrete target is the \texttt{doubleV} PUP family of \cite{comploi2023dynamicAggregates}, which the reference system runs unmodified and which Section~\ref{sec:dsl-limits} shows to fall outside the native fragment on four independent counts. It is worth being explicit about why this is more than a convenience feature. \texttt{doubleV} is expressible today only for the Prolog backend; making it expressible for the native one is what would let the same experiment run on both backends, as every other encoding of this suite does, and \texttt{la\_aux} has already shown that the alternative --- precompiling the missing joins into auxiliary ASP predicates --- reintroduces precisely the grounding product the mechanism exists to avoid. A genuinely two-dimensional rank, restoring the priority field with the safeguards discussed earlier in this chapter, belongs to the same item, since \texttt{doubleV} also annotates a weight that is legitimately zero.

\paragraph{Indexing the foreign store for enumeration-heavy bodies.}
Removing the synchronization cost was future work in an earlier revision of this thesis and is now a completed change, reported in Section~\ref{sec:per-decision-cost}: \texttt{true\_atom/1} and \texttt{false\_atom/1} are registered as non-deterministic foreign predicates that read the assignment from the C++ side, recording a trail change is a hash-map update, and synchronization on \texttt{double-200} fell from $17.5$\,s to $0.63$\,s, level with the native backend. What that change also did was relocate the bottleneck, and the remaining work is where it went.

The residual cost is the snapshot described in Section~\ref{sec:decision-prolog}: a generator call materializes the matching relation before it starts yielding. One snapshot per call is affordable when a rule body calls a generator a handful of times, as on PUP; it is not when the body joins several relations and yields over a thousand candidates per decision, as on HRP, where the per-call cost rose from $9.19$\,ms to $71.9$\,ms and \texttt{lc} lost an instance, \texttt{house-14}. Indexing by signature already keeps the enumeration proportional to the relation rather than to the trail, so what remains is the copy itself. Enumerating the store in place, with a cursor that survives across \texttt{PL\_retry}, would remove it, and is the natural completion of the change --- after which the Prolog backend should be uniformly cheaper than the design it replaced, rather than cheaper on two families and dearer on one.

\paragraph{Instances and benchmarks.}
Separating the \texttt{ga} and \texttt{la} frontiers on PUP requires \texttt{double} instances larger than those published, around $N = 220$ and $N = 240$ on the extrapolation of Section~\ref{sec:pup}. Since the ten reference files follow a regular construction, this means writing a generator and validating it by regenerating $N = 20 \ldots 200$ and diffing byte for byte against the published set before trusting anything it produces beyond it; runs at that size also want a raised time limit, so that a failure of \texttt{ga} is unambiguously a memory failure. Independently of size, multiple instances per size would put variance bars behind the frontier claims, which the present one-run-per-cell protocol cannot support. Beyond PUP, configuration, scheduling, and packing domains are natural candidates for broadening the evaluation, precisely because their useful heuristics are dynamic. For HRP specifically, the instance set released with \cite{semmelrock2026evolvingcode} is the obvious next step: it targets the same configuration domain and would extend the one family of this suite whose sizes are all easy for clasp, which is exactly the limitation recorded above.

\paragraph{Tooling and integration.}
A verified converter from annotated \texttt{\#heuristic} directives to lazy facts, with the flattening rule of Section~\ref{sec:flattening} applied mechanically rather than by hand, would remove the class of error the audit episode discussed earlier in this chapter exemplifies. The mechanism should also become configurable rather than always-on, and the backend should expose its per-decision statistics through clingo's statistics interface instead of stderr parsing. A parallel-solving synchronization layer and a tie-break policy that cooperates with the solver's own scores are the natural solver-side improvements.

\section{Conclusions}

This thesis presented a lazy evaluation mechanism for declarative domain-specific heuristics in clingo, instantiated in two backends over a common propagator design: a native C++ backend interpreting compact \texttt{\_\_heuristic} facts with incremental dynamic aggregates, and an in-process SWI-Prolog backend evaluating heuristic rules as logic programs against the solver trail, which it reads in place through foreign predicates. The mechanism supports two explicit readings of partial information, the clingo-like and the Alpha-like semantics, and reproduces clingo's weight, priority, and modifier behaviour where that reading applies, including the systematic flattening of Alpha's global priority levels into clingo weights.

The full campaign on BSP, PUP, and HRP ran under uniform limits of 600 seconds and 30~GB. It reduces the ground heuristic component from $\Theta(n^3)$ directives to two constant facts, with the corresponding saving in real memory wherever the heuristic grounding matters. On PUP the lazy variant reaches more than twice the solvable frontier of the unguided baseline, and costs half the time and half the memory of the best ground heuristic at every size both solve. On BSP, where the base program owns the frontier, it removes the search cost outright --- $n$ choices and zero conflicts at every size, against tens of thousands for every alternative --- and reaches the same wall as the baseline having spent $0.04$ seconds of solving where the baseline spends $45$. Restating the BSP balance constraint compactly, and running the whole matrix again on it, then separates the two halves of that result: the search advantage disappears, because the reformulated problem needs no guidance, while the representation advantage stands alone and unconfounded at eight million ground rules against six hundred.

The campaign also quantifies the price, on both backends: a per-decision runtime cost that ground-and-solve does not pay. It factorizes into consultation count times consultation cost, ranges from a twentieth of a percent of a run to the whole of it, and differs between the compiled and the interpreted evaluator by two to six orders of magnitude per consultation. It vanishes when the heuristic guides well and dominates when it does not. Within the interpreted backend the two terms proved separable in practice as well as on paper: replacing the mirrored trail with foreign predicates that read the solver's assignment directly removed synchronization as a cost centre altogether, bringing it level with the compiled backend, and moved the residual cost onto rule bodies that enumerate many candidates. The suite validates that all variants and both backends preserve the answer sets, and guards against silent misconfiguration.

Measured against Alpha itself, running the authors' encodings, the lazy propagator reproduces the reference guidance and closes most of the distance to a native lazy-grounding solver: within a factor of $4.5$ in time and at indistinguishable peak memory on PUP, where a ground simulation of the same heuristic is instead nearly nine times slower than the reference and uses almost twice its memory. The residual gap consists entirely of the base program that clingo still grounds.

The final result is a working, documented, and measurable prototype together with a methodology: two orthogonal axes that keep representation and semantics separate, a flattening rule that makes Alpha-style encodings meaningful under clingo's reading of priorities, and an experimental discipline that states what each number measures. It does not solve the grounding bottleneck in general, but it provides a concrete, verifiable path for bringing dynamic declarative heuristics into clingo without always paying the cost of their full ground materialization.
